\textbf{Phần 1. Trắc nghiệm lựa chọn}
\setcounterpageref{ex}{0}

\begin{ex} %[5H8V45]
Một khối gỗ hình hộp chữ nhật được ghép từ $36$ khối lập phương nhỏ cạnh $1\text{ cm}$. Người ta cưa tách khối gỗ đó thành hai phần: phần thứ nhất gồm $15$ khối lập phương nhỏ. Thể tích của phần thứ hai bằng bao nhiêu khối lập phương nhỏ cạnh $1\text{ cm}$?
\choice
{$18\text{ khối}$}
{\True $21\text{ khối}$}
{$24\text{ khối}$}
{$26\text{ khối}$}
\loigiai{
Thể tích phần thứ hai gồm số khối lập phương nhỏ là:
$36 - 15 = 21\text{ (khối)}$.\\
\textbf{Đáp án của câu hỏi là B}
}
\end{ex}

\begin{ex} %[5H8V45]
Hình khối A được tạo bởi $24$ hình lập phương nhỏ cạnh $1\text{ cm}$. Hình khối B có thể tích bằng $\dfrac{3}{4}$ thể tích hình khối A. Thể tích hình khối B gồm bao nhiêu hình lập phương nhỏ cạnh $1\text{ cm}$?
\choice
{$16\text{ khối}$}
{\True $18\text{ khối}$}
{$20\text{ khối}$}
{$32\text{ khối}$}
\loigiai{
Thể tích hình khối B gồm số hình lập phương nhỏ là:
$24 \times \dfrac{3}{4} = 18\text{ (khối)}$.\\
\textbf{Đáp án của câu hỏi là B}
}
\end{ex}

\begin{ex} %[5H8V45]
Người ta xếp các khối lập phương nhỏ cạnh $1\text{ cm}$ thành một hình hộp chữ nhật có chiều dài $5\text{ cm}$, chiều rộng $4\text{ cm}$ và chiều cao $3\text{ cm}$. Tổng số khối lập phương nhỏ cần dùng là
\choice
{$30\text{ khối}$}
{$45\text{ khối}$}
{\True $60\text{ khối}$}
{$120\text{ khối}$}
\loigiai{
Số khối lập phương nhỏ cần dùng là:
$5 \times 4 \times 3 = 60\text{ (khối)}$.\\
\textbf{Đáp án của câu hỏi là C}
}
\end{ex}

\begin{ex} %[5H8V45]
Một khối lập phương lớn được ghép từ $27$ khối lập phương nhỏ cạnh $1\text{ cm}$. Người ta lấy đi ở mỗi đỉnh của khối lập phương lớn đúng một khối lập phương nhỏ. Thể tích của phần còn lại bằng bao nhiêu khối lập phương nhỏ?
\choice
{$23\text{ khối}$}
{$21\text{ khối}$}
{\True $19\text{ khối}$}
{$17\text{ khối}$}
\loigiai{
Khối lập phương có đúng $8$ đỉnh.\\
Số khối lập phương nhỏ bị lấy đi là $8$ khối.\\
Thể tích phần còn lại gồm số khối lập phương nhỏ là:
$27 - 8 = 19\text{ (khối)}$.\\
\textbf{Đáp án của câu hỏi là C}
}
\end{ex}

\begin{ex} %[5H8V45]
Có hai hình khối: hình X được xếp từ $32$ khối lập phương $1\text{ cm}^3$, hình Y được xếp từ $40$ khối lập phương $1\text{ cm}^3$. Tỉ số phần trăm thể tích giữa hình X và hình Y là
\choice
{$75\%$}
{\True $80\%$}
{$85\%$}
{$125\%$}
\loigiai{
Tỉ số phần trăm thể tích hình X so với hình Y là:
$32 : 40 = 0{,}8 = 80\%$.\\
\textbf{Đáp án của câu hỏi là B}
}
\end{ex}

\begin{ex} %[5H8V45]
Để xếp thành một khối lập phương có cạnh dài $3\text{ cm}$ từ các khối lập phương nhỏ cạnh $1\text{ cm}$, người ta cần dùng bao nhiêu khối lập phương nhỏ?
\choice
{$9\text{ khối}$}
{$18\text{ khối}$}
{\True $27\text{ khối}$}
{$36\text{ khối}$}
\loigiai{
Số khối lập phương nhỏ cần dùng là:
$3 \times 3 \times 3 = 27\text{ (khối)}$.\\
\textbf{Đáp án của câu hỏi là C}
}
\end{ex}

\begin{ex} %[5H8V45]
Bạn An có $45$ khối lập phương nhỏ cạnh $1\text{ cm}$. An dùng các khối đó xếp thành một khối lập phương lớn nhất có thể. Hỏi An có thể xếp được khối lập phương có cạnh dài bao nhiêu và còn thừa bao nhiêu khối lập phương nhỏ?
\choice
{Cạnh $4\text{ cm}$ và thừa $5$ khối}
{\True Cạnh $3\text{ cm}$ và thừa $18$ khối}
{Cạnh $3\text{ cm}$ và thừa $20$ khối}
{Cạnh $2\text{ cm}$ và thừa $37$ khối}
\loigiai{
Khối lập phương có cạnh $3\text{ cm}$ cần: $3 \times 3 \times 3 = 27\text{ (khối)}$.\\
Khối lập phương có cạnh $4\text{ cm}$ cần: $4 \times 4 \times 4 = 64\text{ (khối)} > 45$.\\
Vậy khối lập phương lớn nhất có cạnh $3\text{ cm}$ và còn thừa: $45 - 27 = 18\text{ (khối)}$.\\
\textbf{Đáp án của câu hỏi là B}
}
\end{ex}

\begin{ex} %[5H8V45]
Thể tích hình khối A lớn hơn thể tích hình khối B là $8\text{ cm}^3$. Biết thể tích hình A bằng $\dfrac{5}{3}$ thể tích hình B. Thể tích của hình khối B là
\choice
{$8\text{ cm}^3$}
{\True $12\text{ cm}^3$}
{$16\text{ cm}^3$}
{$20\text{ cm}^3$}
\loigiai{
Hiệu số phần bằng nhau là:
$5 - 3 = 2\text{ (phần)}$.\\
Thể tích của hình khối B là:
$8 : 2 \times 3 = 12\text{ (cm}^3\text{)}$.\\
\textbf{Đáp án của câu hỏi là B}
}
\end{ex}

\begin{ex} %[5H8V45]
Một bể cá đang chứa nước. Khi thả chìm hoàn toàn một hòn đá vào trong bể thì mực nước dâng lên cao hơn nhưng không tràn ra ngoài. Thể tích của hòn đá bằng
\choice
{Thể tích toàn bộ nước trong bể cá}
{\True Thể tích của phần nước dâng lên thêm}
{Tổng thể tích của bể cá và lượng nước}
{Diện tích đáy bể cá nhân với chiều cao của bể}
\loigiai{
Thể tích của vật thể chìm hoàn toàn trong nước bằng thể tích của lượng nước dâng lên thêm.\\
\textbf{Đáp án của câu hỏi là B}
}
\end{ex}

\begin{ex} %[5H8V45]
Ghép 4 khối lập phương cạnh $2\text{ cm}$ thành một hàng dài để tạo thành một hình hộp chữ nhật. Thể tích của hình hộp chữ nhật mới tạo thành là
\choice
{$16\text{ cm}^3$}
{$24\text{ cm}^3$}
{\True $32\text{ cm}^3$}
{$64\text{ cm}^3$}
\loigiai{
Thể tích của một khối lập phương cạnh $2\text{ cm}$ là:
$2 \times 2 \times 2 = 8\text{ (cm}^3\text{)}$.\\
Thể tích của 4 khối lập phương ghép lại là:
$8 \times 4 = 32\text{ (cm}^3\text{)}$.\\
\textbf{Đáp án của câu hỏi là C}
}
\end{ex}

\begin{ex} %[5H8V45]
Một hình hộp chữ nhật có kích thước $6\text{ cm} \times 4\text{ cm} \times 5\text{ cm}$ được ghép từ các khối lập phương cạnh $1\text{ cm}$. Người ta sơn đỏ toàn bộ 6 mặt ngoài của hình hộp. Số khối lập phương nhỏ không bị sơn mặt nào là
\choice
{$16\text{ khối}$}
{$20\text{ khối}$}
{\True $24\text{ khối}$}
{$30\text{ khối}$}
\loigiai{
Các khối lập phương không bị sơn mặt nào tạo thành một khối hộp chữ nhật bên trong có kích thước:\\
Chiều dài: $6 - 2 = 4\text{ (cm)}$.\\
Chiều rộng: $4 - 2 = 2\text{ (cm)}$.\\
Chiều cao: $5 - 2 = 3\text{ (cm)}$.\\
Số khối không bị sơn mặt nào là:
$4 \times 2 \times 3 = 24\text{ (khối)}$.\\
\textbf{Đáp án của câu hỏi là C}
}
\end{ex}

\begin{ex} %[5H8V45]
Hai hình khối A và B đều được ghép từ các khối lập phương cạnh $1\text{ cm}$. Khối A gồm 3 tầng, mỗi tầng có 8 khối. Khối B gồm 4 tầng, mỗi tầng có 6 khối. So sánh thể tích hai khối A và B:
\choice
{Thể tích khối A lớn hơn thể tích khối B}
{Thể tích khối B lớn hơn thể tích khối A}
{\True Thể tích hai khối A và B bằng nhau}
{Không so sánh được}
\loigiai{
Thể tích khối A là: $3 \times 8 = 24\text{ (cm}^3\text{)}$.\\
Thể tích khối B là: $4 \times 6 = 24\text{ (cm}^3\text{)}$.\\
Vậy thể tích hai khối bằng nhau.\\
\textbf{Đáp án của câu hỏi là C}
}
\end{ex}

\textbf{Phần 2. Trắc nghiệm đúng sai}
\setcounterpageref{ex}{0}

\begin{ex} %[5H8V45]
Người ta dùng các khối lập phương nhỏ cạnh $1\text{ cm}$ để xếp thành một hình lập phương lớn có cạnh dài $4\text{ cm}$, sau đó sơn màu đỏ lên 6 mặt ngoài của hình lập phương lớn.
\choiceTF
{\True Cần dùng tất cả $64$ khối lập phương nhỏ để xếp thành hình lập phương lớn}
{\True Thể tích của hình lập phương lớn là $64\text{ cm}^3$}
{\True Có đúng $8$ khối lập phương nhỏ nằm ở $8$ đỉnh được sơn đúng $3$ mặt}
{Số khối lập phương nhỏ không bị sơn mặt nào là $16$ khối}
\loigiai{
\begin{itemchoice}
\item Đúng. Số khối cần dùng: $4 \times 4 \times 4 = 64\text{ (khối)}$.
\item Đúng. Thể tích: $64 \times 1\text{ cm}^3 = 64\text{ cm}^3$.
\item Đúng. Tám đỉnh có 8 khối lập phương được sơn 3 mặt.
\item Sai. Số khối ở lõi không bị sơn mặt nào là: $(4 - 2) \times (4 - 2) \times (4 - 2) = 2 \times 2 \times 2 = 8\text{ (khối)} 
eq 16\text{ khối}$.
\end{itemchoice}
}
\end{ex}

\begin{ex} %[5H8V45]
Một bể kính nuôi cá hình hộp chữ nhật có diện tích đáy bên trong là $250\text{ cm}^2$. Người ta thả một hòn non bộ chìm hoàn toàn vào bể thì thấy mực nước dâng cao thêm $2\text{ cm}$ (nước không tràn ra ngoài).
\choiceTF
{\True Chiều cao mực nước dâng lên thêm trong bể là $2\text{ cm}$}
{\True Thể tích của hòn non bộ đúng bằng thể tích của lượng nước dâng lên thêm}
{\True Thể tích của hòn non bộ là $500\text{ cm}^3$}
{Nếu thả thêm một viên đá có thể tích $300\text{ cm}^3$ thì mực nước sẽ dâng thêm đúng $1\text{ cm}$}
\loigiai{
\begin{itemchoice}
\item Đúng. Mực nước dâng thêm $2\text{ cm}$.
\item Đúng. Theo định luật chiếm chỗ của chất lỏng.
\item Đúng. Thể tích non bộ: $250 \times 2 = 500\text{ (cm}^3\text{)}$.
\item Sai. Mực nước dâng thêm là: $300 : 250 = 1{,}2\text{ (cm)} 
eq 1\text{ cm}$.
\end{itemchoice}
}
\end{ex}

\textbf{Phần 3. Trắc nghiệm trả lời ngắn}
\setcounterpageref{ex}{0}

\begin{ex} %[5H8V45]
Cần bao nhiêu khối lập phương nhỏ cạnh $1\text{ cm}$ để xếp thành một hình hộp chữ nhật có chiều dài $6\text{ cm}$, chiều rộng $5\text{ cm}$ và chiều cao $4\text{ cm}$?
\shortans{120}
\loigiai{
Số khối lập phương nhỏ cần dùng là:
$6 \times 5 \times 4 = 120\text{ (khối)}$.
}
\end{ex}

\begin{ex} %[5H8V45]
Một khối lập phương lớn được ghép từ $64$ khối lập phương nhỏ cạnh $1\text{ cm}$. Người ta lấy đi ở mỗi đỉnh của khối lập phương lớn một khối lập phương nhỏ. Hỏi phần còn lại gồm bao nhiêu khối lập phương nhỏ?
\shortans{56}
\loigiai{
Khối lập phương có 8 đỉnh nên đã lấy đi 8 khối nhỏ.\\
Số khối còn lại là:
$64 - 8 = 56\text{ (khối)}$.
}
\end{ex}

\begin{ex} %[5H8V45]
Bạn Dũng có $100$ khối lập phương nhỏ cạnh $1\text{ cm}$. Dũng dùng các khối đó để xếp thành một khối lập phương lớn nhất có thể. Hỏi bạn Dũng còn thừa lại bao nhiêu khối lập phương nhỏ?
\shortans{36}
\loigiai{
Khối lập phương lớn nhất có thể xếp được có cạnh là $4\text{ cm}$, cần dùng:
$4 \times 4 \times 4 = 64\text{ (khối)}$.\\
Số khối còn thừa lại là:
$100 - 64 = 36\text{ (khối)}$.
}
\end{ex}

\begin{ex} %[5H8V45]
Một hình hộp chữ nhật được xếp từ $72$ khối lập phương nhỏ $1\text{ cm}^3$. Biết rằng một lớp đáy của hình hộp chữ nhật đó gồm $24$ khối lập phương nhỏ. Chiều cao của hình hộp chữ nhật đó là bao nhiêu xăng-ti-mét?
\shortans{3}
\loigiai{
Số lớp khối lập phương (chiều cao) là:
$72 : 24 = 3\text{ (cm)}$.
}
\end{ex}

\begin{ex} %[5H8V45]
Người ta xếp các khối lập phương cạnh $1\text{ cm}$ thành một hình lập phương cạnh $3\text{ cm}$ rồi sơn tất cả 6 mặt ngoài. Hỏi có bao nhiêu khối lập phương nhỏ không được sơn mặt nào?
\shortans{1}
\loigiai{
Số khối lập phương nhỏ ở phần tâm bên trong không bị sơn mặt nào là:
$(3 - 2) \times (3 - 2) \times (3 - 2) = 1 \times 1 \times 1 = 1\text{ (khối)}$.
}
\end{ex}

\begin{ex} %[5H8V45]
Một bình đo thể tích chứa $200\text{ cm}^3$ nước. Sau khi thả chìm một viên bi sắt vào trong bình thì nước dâng lên đến vạch $245\text{ cm}^3$. Thể tích của viên bi sắt đó là bao nhiêu xăng-ti-mét khối?
\shortans{45}
\loigiai{
Thể tích của viên bi sắt là:
$245 - 200 = 45\text{ (cm}^3\text{)}$.
}
\end{ex}
